%% =========================================================================
\section{Introduction}\label{sec:intro}

\draftnote{Full draft (2026-07-12), written in plain language. Structure:
P1 why demolition data matters; P2 prior estimates and the indicator behind
each; P3 the problem (BBR has no clean demolition flag); P4 the gap and the
novelty; P5 objectives; P6 short road map. Related work is folded in here
(P2, P4) rather than in a separate section. Remaining markers: \pending{}
for the one headline number the group confirms from results, and
\refneeded{} for the Danish mandatory-LCA citation. RUNE/NICOLAI: please
review and, if useful, add one or two more international references to P2.}

%% ---- P1: why demolition rates matter ------------------------------------
The building sector has a large environmental footprint. Buildings and
their construction account for about a third of the world's energy use and
of its energy-related carbon emissions, and the materials used to build
them cause roughly a further one-fifth of global emissions on their
own~\cite{unep2025gsr}. When a building is demolished, some of its
materials can be recycled or reused, but in practice recovered materials
cover only a small part of what new construction
needs~\cite{andersen2022adaptation}, so much of the carbon already invested
in the building is still lost and new construction is needed to replace
it -- which is why replacing a building generally carries a higher climate
impact than keeping it~\cite{huuhka2023renovate}. How long buildings last, and how many are torn down
each year, therefore matters for several kinds of building-engineering
work. Whole-building life cycle assessment (LCA) uses an assumed building
service life to decide how many times components are replaced and how the
results are reported over time~\cite{ostergaard2018temporal,iso15686}; this
assumption strongly affects the calculated environmental
impact~\cite{andersen2023lifespan}. In Denmark this is no longer only an
academic concern: since 2023 the building regulations require an LCA for
new buildings, calculated over a fixed 50-year period that is set for
comparability rather than to match any building's real
lifespan~\cite{br18_klima}. Models of the building stock and of material flows
also need to know how much is demolished, because demolition is the flow
that leaves the stock and that supplies materials for
reuse~\cite{negendahl2022parametric,andersen2022adaptation}.

%% ---- P2: prior Danish estimates and the indicator behind each -----------
Published estimates of how much is demolished in Denmark, and of how long
Danish buildings last, differ widely. For Denmark alone, reported building
service lives run from medians near 55~years~\cite{ostergaard2018temporal}
to the 120~years assumed in the standard service-life
tables~\cite{aagaard2013levetider}. A key reason is that each study measures
demolition in a different way, even though they draw on the same source: the
Danish Building and Housing Register (Bygnings- og Boligregistret, BBR). The standard Danish service-life tables assume an
overall demolition rate of about 0.3\,\% of the stock per year, taken from
construction-cost statistics rather than from records of individual
buildings~\cite{aagaard2013levetider}. BUILD (Aalborg University) counts
buildings that have been \emph{retired} (``udg{\aa}et'') from the register
and reports on average about 2.2~million~m\textsuperscript{2} of demolished
floor area per year in 2012--2023, again close to 0.3\,\% of the stock by
area~\cite{jensen2025nedrivning}. Studies of building service life instead
start from records of demolition \emph{cases}: one estimates a median life
of roughly 60~years for single-family houses from 124{,}096
cases~\cite{andersen2023lifespan}, and a more recent one finds an
area-weighted median life of 87~years for housing from 106{,}908
demolitions, identified from a demolition-related update to the building
record~\cite{droob2026servicelife}. An earlier study of the temporal scope
of building LCAs used an administrative BBR extract of buildings
``demolished 2009--2015'' without stating how the extract was
made~\cite{ostergaard2018temporal}. Other studies do not use a register
flag at all, and instead detect demolition indirectly, as a house that is
torn down and rebuilt on the same
plot~\cite{jensen2022enfamiliehuse,blv2023karakteristika}. The same pattern
appears abroad: register-based studies in Finland and in a set of European
and US cities (including Copenhagen) each define a demolition as a building
that is removed from, or no longer survives in, a register, rather than as
a directly observed physical event~\cite{huuhka2021stocks,berglundbrown2025lifetimes}.
In short, each of these approaches turns ``a demolition'' into a different
concrete rule -- a register exit, a demolition-coded update, a compiled list
of demolition cases, or a rebuild pattern -- and they also differ in which
of BBR's several floor-area figures they add up.

%% ---- P3: the problem — BBR has no clean demolition flag -----------------
This variety is not careless; it reflects a real limitation of the data. A
register like BBR records administrative objects and processes, not
directly observed physical events, so a change in the register is not the
same as a change on the ground~\cite{mit2021admindata,un2019registers}. BBR
was designed for administration, and property owners are responsible for
reporting changes. The register does contain a purpose-built field that
would mark a completed physical demolition, but that field is not included
in the public data distributed through the national data
portal~\cite{bbr_instruks,datafordeler_grunddatamodel}. What the public data
does contain is a general ``historical'' status that a building also reaches
for other reasons, such as error corrections, mergers, and
re-registrations. There is therefore no single field in the distributed
data that cleanly and only marks a demolition. Asked directly how to tell,
the register authority itself pointed us not to a definitive field but to
the building's lifecycle status -- a proxy. The quality of BBR has also
been found to be uneven: a national audit reported that the extent and
significance of the errors were not even known, and that quality varied
considerably between municipalities~\cite{rigsrevisionen_bbr}, and later
coverage estimated that a large share of single-family houses were recorded
with incorrect measurements~\cite{dr2023bbr}. Faced with this, every
researcher has to build their own \emph{indicator}: an observable register
event, chosen to stand in for a demolition. Because the register offers no
agreed choice, different practitioners choose differently, and this is a
major reason their published numbers do not agree.

%% ---- P4: the gap and the novelty ----------------------------------------
What has not been studied is how much this single choice actually matters.
Earlier sensitivity analyses in building-stock research have varied other
things: the assumed shape of the building-lifespan distribution inside a
material-flow model~\cite{miatto2017buildinglifespan}, or the data
\emph{source} the demolitions are drawn from~\cite{huuhka2021stocks}. To the
best of our knowledge, no study has instead held everything else
fixed -- the same register, the same period, the same calculation -- and
changed only the demolition indicator, in order to measure how much the
resulting national demolition figures move. The novelty of this work is to
make that comparison directly. We treat the choice of indicator as the one
thing we vary, and we ask: how sensitive are Denmark's demolition figures to
how a demolition is defined in BBR? Alongside this main question we ask
several smaller ones: how much do the indicators agree on which buildings
they select; which building types are most affected; what difference excluding
discontinued building-use codes makes; how much the choice
of floor-area figure changes the reported area; and how our transparent
indicators compare with the compiled national extract the field has relied
on.

%% ---- P5: objectives ------------------------------------------------------
Accordingly, this study pursues three objectives. First, to define a set of
demolition indicators directly from a complete BBR extract -- including the
full history of demolished buildings -- so that every step is documented and
can be reproduced, in contrast to the pre-made and undocumented extracts
used in earlier work~\cite{kurvinen2021modeling}. Second, to run the same
demolition-rate calculation once for each indicator, holding all other
choices fixed, and to report how far the results spread. Third, to compare
these results with published Danish figures and with the compiled national
extract, so the practical consequences of the choice are clear. We find that
the choice is not a minor technicality: across indicators the number of
buildings counted as demolished differs by about 4.4-fold, and the reported
demolished floor area depends further on which area figure is used and on
how complete that figure is.

%% ---- P6: road map (optional; keep short) --------------------------------
The remainder of this paper is organised as follows.
Section~\ref{sec:bbr} describes BBR and how a demolition appears in it.
Section~\ref{sec:indicators} defines the demolition indicators we compare.
Section~\ref{sec:methods} explains the data and the shared calculation.
Section~\ref{sec:results} reports the results, and
Sections~\ref{sec:discussion} and~\ref{sec:conclusion} discuss them and
conclude.
